\documentclass[11pt,a4paper]{article}
\usepackage[utf8]{inputenc}
\usepackage[T1]{fontenc}
\IfFileExists{french.ldf}{\usepackage[french]{babel}}{\usepackage[english]{babel}}
\usepackage{amsmath,amssymb,amsthm}
\usepackage{geometry}\geometry{margin=2.3cm}
\usepackage{listings}
\usepackage{xcolor}
\usepackage{enumitem}
\usepackage{fancyhdr}
\usepackage{lmodern}
\usepackage{titlesec}
\usepackage{tikz}
\usetikzlibrary{arrows.meta,positioning}
\usepackage{hyperref}
\hypersetup{colorlinks=true, urlcolor=[RGB]{2,101,115}, linkcolor=black}

\definecolor{codebg}{RGB}{247,249,251}
\definecolor{kw}{RGB}{175,45,45}
\definecolor{cmt}{RGB}{110,120,130}
\definecolor{str}{RGB}{20,110,60}
\definecolor{cnc}{RGB}{2,101,115}
\lstset{
  language=Python, backgroundcolor=\color{codebg}, basicstyle=\ttfamily\small,
  keywordstyle=\color{kw}\bfseries, commentstyle=\color{cmt}\itshape, stringstyle=\color{str},
  numbers=left, numberstyle=\tiny\color{cmt}, numbersep=8pt, frame=single, rulecolor=\color{gray!40},
  showstringspaces=false, breaklines=true, columns=fullflexible, extendedchars=true, inputencoding=utf8,
  literate={é}{{\'e}}1 {è}{{\`e}}1 {à}{{\`a}}1 {ê}{{\^e}}1 {î}{{\^i}}1 {ç}{{\c c}}1 {ù}{{\`u}}1 {ô}{{\^o}}1 {û}{{\^u}}1 {'}{{'}}1
}

\newcounter{qcnt}
\newcommand{\Q}{\medskip\par\stepcounter{qcnt}\noindent\textbf{Question~\theqcnt.}~}
\newcommand{\code}[1]{\lstinline!#1!}
\newcommand{\hh}[1]{h(#1)}

\newtheorem{lemme}{Lemme}
\newtheorem{theoreme}{Théorème}
\theoremstyle{definition}
\newtheorem{defi}{Définition}

\pagestyle{fancy}\fancyhf{}
\lhead{\small Préparation au CNC --- Informatique MP}
\rhead{\small Protocoles TCP/IP et HTTP/HTTPS}
\cfoot{\thepage}
\renewcommand{\headrulewidth}{0.4pt}
\titleformat*{\section}{\large\bfseries\scshape\color{cnc}}
\titleformat*{\subsection}{\bfseries\color{cnc}}
\usepackage{longtable}
\newcommand{\rfc}[1]{\href{https://www.rfc-editor.org/rfc/rfc#1}{\textsc{rfc}~#1}}

\rhead{\small TCP/IP et HTTP/HTTPS --- Corrigé}
\renewcommand{\Q}{\medskip\par\stepcounter{qcnt}\noindent\textbf{\color{cnc}Question~\theqcnt.}~}
\lstdefinestyle{sql}{language=SQL, morekeywords={SUBSTR,ROUND,GROUP,ORDER,LIMIT,HAVING,CASE,WHEN,THEN,ELSE,END,EXISTS,DISTINCT}}
\begin{document}
\begin{center}
{\large\bfseries Préparation au Concours National Commun --- Informatique MP}\\[8pt]
\rule{0.9\linewidth}{0.5pt}\\[6pt]
{\Large\bfseries Corrigé détaillé}\\[4pt]
{\bfseries Filière MP --- Durée : 4 heures}\\[6pt]
{\itshape Voyage d'un paquet : de l'adresse IP au cadenas du navigateur}\\[4pt]
{\itshape Les protocoles TCP/IP et HTTP/HTTPS}\\[6pt]
\rule{0.9\linewidth}{0.5pt}\\[8pt]
{\small Auteur~: \href{https://orcid.org/0000-0002-6108-7176}{\textbf{Pr Agrégé El Hadiq Zouhair}}}\\[3pt]
{\small \href{https://cpgeacademy.org}{\textbf{cpgeacademy.org}} \quad---\quad Tous droits réservés \textcopyright\ cpgeacademy.org}
\end{center}
\medskip
\begin{center}
\fcolorbox{gray}{gray!5}{\begin{minipage}{0.93\linewidth}\small
\textbf{Avis important.} Ce document est une \textbf{initiative personnelle} du professeur,
à but strictement pédagogique. Il ne s'agit \textbf{pas} d'un sujet officiel du Concours
National Commun.
\smallskip

\textbf{Note de vérification.} L'intégralité du code Python de ce corrigé a été exécutée~;
les fonctions ont été confrontées à des implantations naïves de référence sur plusieurs
milliers de jeux d'essai aléatoires (arbre de préfixes, fusion d'intervalles, cache LRU,
programmation dynamique), et les requêtes SQL ont été exécutées sur une base réelle. Les
exemples numériques des questions 16 et 17 sont ceux des \rfc{1624} et \rfc{1071}.
\smallskip

\textbf{Pour aller plus loin.} \href{https://cpgeacademy.org}{\textbf{cpgeacademy.org}}.
\end{minipage}}
\end{center}
\medskip

%=======================================================================
\section*{Partie I --- Adressage IPv4 et agrégation CIDR}
%=======================================================================

\Q Il s'agit d'une écriture en base $256$ sur exactement quatre chiffres.
\begin{lstlisting}
def ip_vers_entier(s):
    x = 0
    for c in s.split('.'):
        x = x * 256 + int(c)
    return x

def entier_vers_ip(x):
    o = []
    for _ in range(4):
        o.append(str(x % 256))
        x = x // 256
    return '.'.join(reversed(o))
\end{lstlisting}
Les deux boucles font exactement quatre tours~: la complexité est $\Theta(1)$.
(\code{ip_vers_entier} applique le schéma de Horner en base $256$~;
\code{entier_vers_ip} extrait les chiffres du poids faible vers le poids fort, d'où le
\code{reversed}.)

\Q
\begin{lstlisting}
def masque(n):
    if n == 0:
        return 0
    return ((1 << n) - 1) << (32 - n)
\end{lstlisting}
\textbf{Démonstration.} Pour $n\ge 1$, $2^{n}-1$ est l'entier dont les $n$ bits de poids
faible valent $1$. Le décaler de $32-n$ positions vers la gauche revient à le multiplier par
$2^{32-n}$~: on obtient $(2^{n}-1)\,2^{32-n}=2^{32}-2^{32-n}$, entier dont les $n$ bits de
poids fort valent $1$ et les $32-n$ autres $0$. Le cas $n=0$ est traité à part car
$\texttt{1 << 0} - 1 = 0$ conviendrait, mais un décalage de $32$ n'est pas défini sur un mot
de $32$ bits dans la plupart des langages~; on garde donc le test pour la portabilité.
\hfill$\square$

\Q
\begin{lstlisting}
def adresse_reseau(a, n):
    return a & masque(n)

def diffusion(a, n):
    return a | (0xFFFFFFFF ^ masque(n))
\end{lstlisting}
\textbf{Démonstration.} Notons $r=\texttt{adresse\_reseau}(a,n)$. Par définition du et bit à
bit, $r$ a les mêmes $n$ bits de poids fort que $a$ et ses $32-n$ bits de poids faible nuls.
De même, $\texttt{diffusion}(a,n)$ a les mêmes $n$ bits de poids fort que $a$ et ses $32-n$
bits de poids faible tous égaux à $1$, c'est-à-dire vaut $2^{32-n}-1$ sur cette partie. Comme
$r$ vaut $0$ sur ces mêmes bits et que les deux nombres coïncident ailleurs,
\[ \texttt{diffusion}(a,n)-r \;=\; 2^{32-n}-1 . \]
Ce sont bien le minimum et le maximum de $\mathcal{A}(p/n)$~: écrire $a=r+u$ avec
$0\le u\le 2^{32-n}-1$ établit que $\mathcal{A}(p/n)=[\![r,\ r+2^{32-n}-1]\!]$, ensemble de
cardinal $2^{32-n}$. \hfill$\square$

\Q Le préfixe contient $2^{32-n}$ adresses~; deux sont réservées (réseau et diffusion), d'où
\[ \#\{\text{adresses attribuables}\} \;=\; 2^{32-n}-2 \qquad (n\le 30). \]
Pour $n=31$ la formule donnerait $0$~: un tel préfixe ne contient que deux adresses, toutes
deux « réservées ». Cette situation est traitée à part (\rfc{3021})~: sur une liaison
point-à-point il n'y a pas besoin d'adresse de diffusion, et les deux adresses sont
attribuées aux deux extrémités. Pour $n=32$ le préfixe se réduit à une adresse unique~: c'est
une \emph{route d'hôte}. Ces deux cas montrent que la formule $2^{32-n}-2$ est une convention
d'ingénierie et non un théorème.

\Q
\begin{lstlisting}
def dans_prefixe(a, p, n):
    if n == 0:
        return True
    return (a ^ p) >> (32 - n) == 0
\end{lstlisting}
\textbf{Démonstration.} Notons $a_{31},\dots,a_0$ et $p_{31},\dots,p_0$ les bits de $a$ et
$p$. Le ou exclusif est bit à bit~: le bit de rang $i$ de $a\oplus p$ vaut $0$ si et
seulement si $a_i=p_i$. Le décalage $\texttt{>>}(32-n)$ ne conserve que les bits de rang
$\ge 32-n$, c'est-à-dire les $n$ bits de poids fort. Donc
\[ (a\oplus p)\,\texttt{>>}\,(32-n)=0 \iff \forall i\ge 32-n,\ a_i=p_i \iff
a\in\mathcal{A}(p/n). \]
\hfill$\square$ Complexité $\Theta(1)$~; on n'a besoin ni de calculer le masque ni de faire
deux et logiques.

\Q \textbf{(Agrégation.)}

$(\Leftarrow)$ Supposons que $p$ et $q$ ne diffèrent que par le bit de rang $32-n$~; quitte à
les échanger, $p<q$ et $q=p+2^{32-n}$. D'après la question 3,
$\mathcal{A}(p/n)=[\![p,\,p+2^{32-n}-1]\!]$ et $\mathcal{A}(q/n)=[\![p+2^{32-n},\,
p+2^{32-n+1}-1]\!]$. Ces deux intervalles sont contigus, leur réunion est
$[\![p,\,p+2^{33-n}-1]\!]$. Or $p$ a ses $32-n$ bits de poids faible nuls \emph{et} son bit
de rang $32-n$ nul (puisque $q=p+2^{32-n}$ n'a pas provoqué de retenue), donc $p$ a ses
$33-n$ bits de poids faible nuls~: $p/(n-1)$ est un préfixe valide et
$\mathcal{A}(p/(n-1))=[\![p,\,p+2^{33-n}-1]\!]$. On prend $r=p$.

$(\Rightarrow)$ Réciproquement, supposons $\mathcal{A}(r/(n-1))=\mathcal{A}(p/n)\cup
\mathcal{A}(q/n)$. Le membre de gauche a pour cardinal $2^{33-n}$ et le membre de droite au
plus $2\cdot 2^{32-n}=2^{33-n}$, avec égalité si et seulement si les deux préfixes sont
disjoints. Ils le sont donc, et leur réunion est un intervalle~: ils sont contigus, donc
$q=p+2^{32-n}$ (à échange près), ce qui signifie exactement que $p$ et $q$ ne diffèrent que
par le bit de rang $32-n$. \hfill$\square$

\textbf{Intérêt.} Un opérateur qui possède deux préfixes adjacents n'en annonce qu'un seul,
deux fois plus grand~: c'est l'agrégation CIDR, sans laquelle la table de routage globale
d'Internet, aujourd'hui de l'ordre de $10^{6}$ entrées, en compterait plusieurs dizaines de
millions.

\Q \textbf{(Famille laminaire.)} Supposons $n\le m$ (sinon on échange les rôles) et
$\mathcal{A}(p/n)\cap\mathcal{A}(q/m)\ne\varnothing$~; soit $a$ dans l'intersection. Alors
$a$ coïncide avec $p$ sur ses $n$ bits de poids fort, et avec $q$ sur ses $m\ge n$ bits de
poids fort. Donc $p$ et $q$ coïncident sur leurs $n$ bits de poids fort. Soit alors
$b\in\mathcal{A}(q/m)$~: $b$ coïncide avec $q$ sur $m\ge n$ bits, donc avec $q$ — donc avec
$p$ — sur $n$ bits, c'est-à-dire $b\in\mathcal{A}(p/n)$. Ainsi
$\mathcal{A}(q/m)\subseteq\mathcal{A}(p/n)$.

Les trois cas sont donc exhaustifs~; ils sont mutuellement exclusifs dès que les deux
préfixes sont distincts, puisque les ensembles sont non vides. \hfill$\square$

\emph{Autrement dit, l'ensemble des préfixes CIDR ordonné par inclusion est une forêt~: c'est
exactement ce qui permet la structure d'arbre de la question 20.}

\Q
\begin{enumerate}[label=(\alph*),leftmargin=2em,itemsep=2pt]
\item $\texttt{masque}(12)=2^{32}-2^{20}$. En base $256$~: les $12$ premiers bits sont à $1$,
soit $\texttt{11111111.1111}0000.0.0$, c'est-à-dire
$\boxed{\texttt{255.240.0.0}}$ (car $\texttt{11110000}_2=240$).
\item $\texttt{172.16.35.7}$ et le masque $\texttt{255.240.0.0}$~: $172\ \&\ 255=172$,
$16\ \&\ 240 = 16$ (car $16=\texttt{00010000}_2$), $35\ \&\ 0=0$, $7\ \&\ 0=0$. L'adresse de
réseau est $\boxed{\texttt{172.16.0.0}}$.
\item $\texttt{masque}(26)=\texttt{255.255.255.192}$ (car $192=\texttt{11000000}_2$). Or
$130=\texttt{10000010}_2$, donc $130\ \&\ 192 = 128$~: le réseau est
$\boxed{\texttt{192.168.1.128}}$. La diffusion s'obtient en mettant à $1$ les $6$ bits de
poids faible~: $128+63=191$, soit $\boxed{\texttt{192.168.1.191}}$. Le nombre d'adresses
utilisables est $2^{6}-2=\boxed{62}$.
\end{enumerate}

%=======================================================================
\section*{Partie II --- La somme de contrôle Internet}
%=======================================================================

\Q
\begin{lstlisting}
def add16(a, b):
    s = a + b
    return (s & 0xFFFF) + (s >> 16)
\end{lstlisting}
\textbf{Démonstration.} Soient $a,b\in[\![0,N-1]\!]$ avec $N=2^{16}$. Alors
$s=a+b\le 2N-2$, donc $\lfloor s/N\rfloor\in\{0,1\}$~; l'expression \code{(s >> 16)} vaut
donc $0$ ou $1$, et le résultat est $(s\bmod N)+\lfloor s/N\rfloor$.

\emph{Cas $s<N$~:} le résultat vaut $s\le N-1$.
\emph{Cas $s\ge N$~:} le résultat vaut $s-N+1$, et $N\le s\le 2N-2$ donne
$1\le s-N+1\le N-1$.

Dans les deux cas le résultat est dans $[\![0,N-1]\!]$~: \textbf{la réinjection de la retenue
ne peut jamais provoquer une seconde retenue}. En effet, elle ne pourrait le faire que si
$s\bmod N=N-1$ et $\lfloor s/N\rfloor=1$, c'est-à-dire $s=2N-1$, impossible puisque
$s\le 2N-2$. \hfill$\square$

\Q Reprenons les deux cas. Si $s<N$, le résultat est $s=a+b$. Si $s\ge N$, le résultat est
$s-(N-1)$. Dans les deux cas
\[ a\oplus' b \;\equiv\; a+b \pmod{N-1}. \]
La commutativité est immédiate ($a+b=b+a$). Pour l'associativité, remarquons que
$(a\oplus' b)\oplus' c$ et $a\oplus'(b\oplus' c)$ sont deux éléments de $[\![0,N-1]\!]$
congrus à $a+b+c$ modulo $N-1$~; deux tels éléments diffèrent d'un multiple de $N-1$, donc
sont égaux \emph{ou} valent $\{0,N-1\}$. Un examen des cas montre que $\oplus'$ ne produit
$0$ que si les deux opérandes sont nuls~: si $s\ge N$ le résultat est $\ge 1$, et si $s<N$ le
résultat est $s$, nul seulement si $a=b=0$. Comme $a\oplus'0=a$, les deux parenthésages
donnent $0$ exactement dans le même cas ($a=b=c=0$), et sont donc toujours égaux.
\hfill$\square$

$0$ est neutre~: $a\oplus'0=a$. Enfin $x\oplus'\overline{x}=x+(N-1-x)=N-1=\texttt{0xFFFF}$.

\textbf{Les deux zéros.} $\texttt{0x0000}$ et $\texttt{0xFFFF}$ sont tous deux congrus à $0$
modulo $N-1$~: ce sont les deux représentations de zéro, notées $+0$ et $-0$. L'opération
$\oplus'$ ne produit jamais $+0$ à partir d'opérandes non tous nuls~: elle « préfère »
$-0$. Cette dissymétrie est \emph{exactement} la source du bogue de la \rfc{1141} étudié à la
question 16.

\Q \textbf{(a) Découpage en mots.}
\begin{lstlisting}
def mots16(octets):
    o = list(octets)                     # copie : on ne modifie pas l'argument
    if len(o) % 2 == 1:
        o.append(0)                      # bourrage a droite
    return [o[i]*256 + o[i+1] for i in range(0, len(o), 2)]
\end{lstlisting}
L'octet d'indice pair est celui de poids fort~: c'est la convention « gros-boutienne » de
l'ordre réseau. Sur l'exemple demandé,
$\texttt{mots16([0x45, 0x00, 0x00, 0x73])} = [\texttt{0x4500},\ \texttt{0x0073}]$.

\emph{On prendra garde à recopier la liste~: sans le \code{list(octets)}, le bourrage
modifierait la liste de l'appelant, effet de bord d'autant plus pernicieux qu'il ne se
manifeste que pour les paquets de longueur impaire.}

\textbf{(b) Somme.}
\begin{lstlisting}
def somme_un(mots):
    s = 0
    for m in mots:
        s = add16(s, m)
    return s
\end{lstlisting}
L'accumulateur est initialisé à $0$, élément neutre de $\oplus'$ (question 10).

\textbf{(c) Complément.}
\begin{lstlisting}
def checksum(octets):
    return somme_un(mots16(octets)) ^ 0xFFFF
\end{lstlisting}
Sur l'exemple~: $\texttt{0x4500}\oplus'\texttt{0x0073}=\texttt{0x4573}$ (aucune retenue),
d'où $\texttt{checksum}=\overline{\texttt{0x4573}}=\texttt{0xFFFF}-\texttt{0x4573}
=\boxed{\texttt{0xBA8C}}$.

\textbf{Complexités.} Notons $k$ le nombre d'octets et $\ell=\lceil k/2\rceil$ le nombre de
mots.
\begin{center}\small
\begin{tabular}{l|c|c}
 & temps & espace auxiliaire\\\hline
\code{mots16} & $\Theta(k)$ & $\Theta(k)$ (la copie et la liste de mots)\\
\code{somme_un} & $\Theta(\ell)=\Theta(k)$ & $\Theta(1)$ (un accumulateur)\\
\code{checksum} & $\Theta(k)$ & $\Theta(k)$ (héritée de \code{mots16})
\end{tabular}
\end{center}

\textbf{Version en espace constant.} La liste intermédiaire des mots ne sert à rien~: on peut
les former et les accumuler dans la même boucle.
\begin{lstlisting}
def checksum_direct(octets):
    s = 0
    k = len(octets)
    i = 0
    while i + 1 < k:                     # paires completes
        s = add16(s, octets[i]*256 + octets[i+1])
        i = i + 2
    if i < k:                            # dernier octet isole
        s = add16(s, octets[i]*256)      # bourrage implicite par un octet nul
    return s ^ 0xFFFF
\end{lstlisting}
On obtient $\Theta(k)$ en temps et $\boxed{\Theta(1)}$ en espace auxiliaire, sans modifier
l'argument. C'est exactement la boucle donnée en annexe du \rfc{1071}~; sur un paquet de
$1500$ octets, l'économie de la liste intermédiaire est loin d'être négligeable lorsqu'on
traite des millions de paquets par seconde.

\Q Notons $S$ la somme en complément à un de tous les mots, le champ de contrôle étant à
zéro. La valeur inscrite est $\overline{S}$. Le récepteur calcule la somme du paquet complet,
c'est-à-dire, par associativité et commutativité (question 10)~:
\[ S\oplus'\overline{S} \;=\; N-1 \;=\; \texttt{0xFFFF}. \]
(On a utilisé $x\oplus'\overline{x}=\texttt{0xFFFF}$, établi à la question 10.)
\hfill$\square$
\begin{lstlisting}
def verifie(octets):
    return somme_un(mots16(octets)) == 0xFFFF
\end{lstlisting}
On notera l'élégance du procédé~: le récepteur n'a \emph{pas} à extraire le champ de
contrôle ni à le remettre à zéro. Il fait la même boucle que l'émetteur et compare à une
constante.

\Q \textbf{(Indépendance de l'ordre des octets.)} Écrivons $x=256u+v$ avec
$u,v\in[\![0,255]\!]$~; alors $\sigma(x)=256v+u$.

Il suffit de traiter le cas de deux mots, le cas général s'en déduisant par récurrence
immédiate grâce à l'associativité. Soient $x=256u+v$ et $y=256u'+v'$.

\emph{Cas 1~: $v+v'\le 255$ et $u+u'\le 255$.} Aucune retenue nulle part~:
$x\oplus' y=256(u+u')+(v+v')$ et $\sigma(x)\oplus'\sigma(y)=256(v+v')+(u+u')$, qui est bien
son échangé.

\emph{Cas 2~: retenue de l'octet de poids faible vers celui de poids fort.} Dans le calcul
échangé, cette même retenue passe de l'octet de poids fort vers l'octet de poids faible
\emph{par la réinjection circulaire}. Le point clé, souligné par la \rfc{1071}, est que
l'addition en complément à un sur $16$ bits est \emph{circulaire}~: les retenues vont du bit
$7$ au bit $8$ et du bit $15$ au bit $0$. Échanger systématiquement les octets revient à
tourner les bits de $8$ positions dans ce cycle, ce qui ne modifie pas leur ordre relatif~:
la somme obtenue est la somme échangée. \hfill$\square$

\textbf{Conséquence pratique.} Une machine « petit-boutienne » qui lit en mémoire, mot par
mot, des données stockées en ordre réseau (« gros-boutien ») obtient automatiquement la somme
échangée~; en réécrivant le résultat en mémoire, il est de nouveau échangé, donc correct. La
somme de contrôle se calcule ainsi avec \emph{exactement le même code} quelle que soit
l'architecture — argument décisif en 1988.

\Q Supposons que le paquet reçu diffère du paquet émis par un seul bit, de rang
$r\in[\![0,15]\!]$ dans le mot d'indice $i$. La somme reçue $S'$ vérifie
$S'\equiv S\pm 2^{r}\pmod{N-1}$, où $S=\texttt{0xFFFF}\equiv 0$. Or $1\le 2^{r}\le 2^{15}$,
donc $2^{r}\not\equiv 0 \pmod{N-1}$ (car $0<2^r<N-1$), et de même $-2^r\not\equiv 0$. Donc
$S'\not\equiv 0$, c'est-à-dire $S'\ne\texttt{0xFFFF}$ et $S'\ne\texttt{0x0000}$~: le test de
la question 12 échoue. \hfill$\square$

\emph{Vérification expérimentale~: sur $3\,000$ paquets aléatoires de $20$ octets et une
inversion de bit tirée au hasard, l'erreur a été détectée dans $3\,000$ cas sur $3\,000$.}

\Q L'addition $\oplus'$ étant \textbf{commutative}, permuter deux mots de $16$ bits ne change
pas la somme~: une telle altération n'est \emph{jamais} détectée. Par exemple, le paquet
commençant par $\texttt{12 34 AB CD}$ et celui commençant par $\texttt{AB CD 12 34}$, tout le
reste étant identique, ont la même somme de contrôle.

De même, une erreur de deux bits qui \emph{se compense} passe inaperçue~: ajouter $1$ à un
mot et retrancher $1$ à un autre laisse la somme inchangée. Ainsi $\texttt{12 35 AB CC}$ a la
même somme que $\texttt{12 34 AB CD}$.

\textbf{Comparaison avec un CRC.} Un code polynomial interprète le message comme un polynôme
sur $\mathbb{F}_2$ et en calcule le reste modulo un polynôme générateur~: le résultat dépend
de la \emph{position} de chaque bit, ce qui détruit la symétrie ci-dessus. Un CRC-32 bien
choisi détecte toutes les erreurs de $1$, $2$ ou $3$ bits, toutes les rafales de longueur
$\le 32$ et toutes les erreurs de poids impair. Il est en revanche plus coûteux à calculer et
ne se met pas à jour incrémentalement (question 16). Le choix de la somme en complément à un
pour la couche transport est donc un compromis assumé~: la détection fine est confiée au
CRC de la couche liaison, la somme de contrôle Internet n'assurant qu'un contrôle
\emph{de bout en bout} bon marché.

\Q \textbf{(Mise à jour incrémentale.)}
\begin{lstlisting}
def maj_1141(HC, m, mp):          # HC' = HC +' m +' ~m'      (RFC 1141)
    return add16(add16(HC, m), mp ^ 0xFFFF)

def maj_1624(HC, m, mp):          # HC' = ~( ~HC +' ~m +' m' ) (RFC 1624)
    return add16(add16(HC ^ 0xFFFF, m ^ 0xFFFF), mp) ^ 0xFFFF
\end{lstlisting}
\textbf{(c) Application numérique.} On a $C_{\text{autres}}=\texttt{0xCD7A}$,
$m=\texttt{0x5555}$, $m'=\texttt{0x3285}$. On note d'abord les deux sommes utiles~:
\[ \texttt{0xCD7A}+\texttt{0x5555}=\texttt{0x122CF}\ \longrightarrow\ \texttt{0x22D0}
   \quad\text{(repli)},\qquad
   \texttt{0xCD7A}+\texttt{0x3285}=\texttt{0xFFFF}. \]
On en déduit~:
\[
\begin{array}{r@{\ }c@{\ }l}
HC &=& \overline{\texttt{0x22D0}} \;=\; \texttt{0xDD2F},\\[4pt]
HC'\ \text{(recalcul complet)} &=& \overline{\texttt{0xFFFF}} \;=\;
   \boxed{\texttt{0x0000}},\\[4pt]
HC'\ \text{par la \rfc{1141}} &=& \texttt{0xDD2F}\oplus'\texttt{0x5555}\oplus'
   \overline{\texttt{0x3285}} \;=\; \texttt{0xFFFF} \qquad\textbf{faux},\\[4pt]
HC'\ \text{par la \rfc{1624}} &=& \overline{\ \texttt{0x22D0}\oplus'
   \overline{\texttt{0x5555}}\oplus'\texttt{0x3285}\ } \;=\; \texttt{0x0000}
   \qquad\textbf{correct}.
\end{array}
\]

\textbf{(d) Explication.} La formule de la \rfc{1141} suppose que la complémentation est
\emph{distributive} sur $\oplus'$, c'est-à-dire $\overline{x\oplus' y}=\overline{x}\oplus'
\overline{y}\oplus'\dots$~; or cette propriété tombe en défaut lorsque le résultat est nul,
précisément à cause des deux représentations de zéro (question 10). L'addition $\oplus'$ ne
produit jamais $+0$ à partir d'opérandes non tous nuls~: elle produit $-0=\texttt{0xFFFF}$.
Comme l'en-tête contient toujours au moins un champ non nul, le champ de contrôle ne peut
jamais valoir $\overline{+0}=\texttt{0xFFFF}$~; il peut en revanche valoir
$\overline{-0}=\texttt{0x0000}$. La \rfc{1141} produit donc $\texttt{0xFFFF}$, une valeur
qu'un calcul complet ne peut jamais donner. La \rfc{1624} évite le problème en n'utilisant la
complémentation qu'une seule fois, tout à la fin.

\emph{Remarque.} Le \rfc{1624} note que si le récepteur vérifie comme à la question 12 (somme
du paquet complet $=\texttt{0xFFFF}$), les deux valeurs $\texttt{0x0000}$ et
$\texttt{0xFFFF}$ passent le test~: le bogue n'a d'effet que sur les implantations qui
recalculent la somme et la comparent au champ. C'est un bel exemple de bogue qui ne se
manifeste qu'à l'interopérabilité.

\Q On regroupe les $20$ octets en dix mots de $16$ bits~:
\[ \texttt{4500},\ \texttt{0073},\ \texttt{0000},\ \texttt{4000},\ \texttt{4011},\
   \texttt{0000},\ \texttt{c0a8},\ \texttt{0001},\ \texttt{c0a8},\ \texttt{00c7}. \]
La question 10 autorise à sommer d'abord \emph{sans} se soucier des retenues, puis à les
replier à la fin (c'est la technique dite des « retenues différées » du \rfc{1071}). La somme
entière vaut
\[ 17664+115+0+16384+16401+0+49320+1+49320+199 \;=\; 149\,404 \;=\; \texttt{0x2479C}. \]
On replie la retenue de poids fort~: $\texttt{0x479C}+\texttt{0x2}=\texttt{0x479E}$, valeur
qui tient sur $16$ bits — aucun second repli n'est nécessaire. D'où
\[ \boxed{\texttt{checksum} \;=\; \overline{\texttt{0x479E}} \;=\; \texttt{0xFFFF}-
\texttt{0x479E} \;=\; \texttt{0xB861}} \]
que l'on inscrit dans les octets d'indices $10$ et $11$~: \texttt{b8 61}.

\textbf{Vérification.} En recalculant la somme sur le paquet complet, on ajoute
$\texttt{0xB861}$ à $\texttt{0x479E}$, ce qui donne exactement $\texttt{0xFFFF}$~: le test de
la question 12 est satisfait. (C'est l'exemple classique du \rfc{1071}~; le calcul a été
vérifié en machine.)

%=======================================================================
\section*{Partie III --- Routage : plus long préfixe et plus courts chemins}
%=======================================================================

\Q
\begin{lstlisting}
def plus_long_prefixe(table, a):
    best, bl = None, -1
    for (p, n, saut) in table:
        if dans_prefixe(a, p, n) and n > bl:
            bl, best = n, saut
    return best
\end{lstlisting}
Chaque entrée est examinée une fois, en $O(1)$ (question 5)~: la complexité est
$\Theta(k)$.

\Q Soit $\mathcal{P}$ l'ensemble des préfixes de la table contenant $a$. Il est non vide (il
contient la route par défaut). D'après la question 7, deux éléments de $\mathcal{P}$ sont
comparables pour l'inclusion, puisque leur intersection contient $a$ donc est non vide~:
$\mathcal{P}$ est \textbf{totalement ordonné} par inclusion. Comme
$\mathcal{A}(p/n)\subseteq\mathcal{A}(q/m)\Rightarrow n\ge m$, l'ordre par inclusion coïncide
avec l'ordre inverse des longueurs~: il existe donc un unique élément de longueur maximale,
qui est aussi le plus petit pour l'inclusion. \hfill$\square$

\emph{C'est ce résultat qui donne son sens à la règle « plus long préfixe »~: elle sélectionne
la route la plus spécifique, et cette route est unique.}

\Q On représente un n\oe ud par une liste \code{[fils0, fils1, saut]}.
\begin{lstlisting}
class Trie:
    def __init__(self):
        self.racine = [None, None, None]

    def inserer(self, p, n, saut):
        noeud = self.racine
        for i in range(n):
            b = (p >> (31 - i)) & 1          # i-eme bit en partant du poids fort
            if noeud[b] is None:
                noeud[b] = [None, None, None]
            noeud = noeud[b]
        noeud[2] = saut
\end{lstlisting}
\textbf{Complexités.} Une insertion effectue exactement $n\le 32$ tours en $O(1)$~: elle
coûte $O(n)=O(32)=O(1)$ si l'on considère $32$ comme une constante. En espace, chaque
insertion crée au plus $32$ n\oe uds~: la structure occupe $O(32k)=O(k)$ pour $k$ entrées.
Le facteur $32$ est le prix de la structure, mais il se réduit fortement en pratique par
compression des chaînes de n\oe uds à un seul fils (arbres \emph{Patricia}).

\Q
\begin{lstlisting}
    def chercher(self, a):
        noeud = self.racine
        res = self.racine[2]                 # eventuelle route par defaut (/0)
        for i in range(32):
            b = (a >> (31 - i)) & 1
            if noeud[b] is None:
                break
            noeud = noeud[b]
            if noeud[2] is not None:
                res = noeud[2]
        return res
\end{lstlisting}
\textbf{Invariant.}
\begin{quote}\itshape
Au début de l'itération $i$, \code{noeud} est le n\oe ud de profondeur $i$ correspondant aux
$i$ premiers bits de $a$ (ce n\oe ud existe), et \code{res} est le saut associé au plus long
préfixe de longueur $\le i$ de la table qui contienne $a$ ($\texttt{None}$ s'il n'y en a pas).
\end{quote}
\emph{Initialisation.} Pour $i=0$, \code{noeud} est la racine, qui représente la suite vide,
et \code{res} est le saut du préfixe $/0$ s'il existe.

\emph{Hérédité.} Supposons l'invariant vrai au début de l'itération $i$. Soit $b$ le bit de
rang $31-i$ de $a$. Si $\texttt{noeud[b]}$ est \code{None}, aucun préfixe de la table de
longueur $>i$ ne peut contenir $a$ (par construction de \code{inserer}, un tel préfixe aurait
créé ce fils)~: la boucle s'arrête et \code{res} est déjà la bonne valeur. Sinon on descend~;
le nouveau n\oe ud représente les $i+1$ premiers bits de $a$. S'il porte un saut, c'est que
la table contient le préfixe correspondant, de longueur $i+1$, qui contient $a$ et qui est
plus long que tous les précédents~: on met à jour \code{res}. L'invariant est rétabli pour
$i+1$.

\emph{Conclusion.} À la sortie, \code{res} est le saut du plus long préfixe de la table
contenant $a$, qui est unique d'après la question 19. \hfill$\square$

\emph{Vérification expérimentale~: sur $300$ tables aléatoires et $200$ adresses par table,
soit $60\,000$ recherches, l'arbre et la version exhaustive de la question 18 ont toujours
donné le même résultat.}

\Q La boucle fait au plus $32$ tours en $O(1)$~: la complexité est $\boxed{O(32)=O(1)}$,
\textbf{indépendante du nombre $k$ d'entrées}, contre $\Theta(k)$ pour la question 18.

\textbf{Commentaire.} Un routeur de c\oe ur de réseau porte aujourd'hui une table d'environ
$10^{6}$ préfixes et doit décider du sort de plusieurs dizaines de millions de paquets par
seconde. Le parcours exhaustif demanderait $10^{6}$ tests par paquet — impossible. La
recherche dans l'arbre en demande $32$ au plus, soit un gain de plus de quatre ordres de
grandeur. C'est exactement l'écart entre un algorithme linéaire et un algorithme
logarithmique en la taille de l'espace d'adressage. En pratique les routeurs vont plus loin
encore (arbres multi-bits, mémoires associatives), mais le principe reste celui-ci.

\Q
\begin{lstlisting}
import heapq

def dijkstra(V, adj, s):
    INF = float('inf')
    d = [INF] * V
    d[s] = 0
    vus = [False] * V
    tas = [(0, s)]
    while tas:
        du, u = heapq.heappop(tas)
        if vus[u]:
            continue                          # entree perimee
        vus[u] = True
        for (v, w) in adj[u]:
            if d[u] + w < d[v]:
                d[v] = d[u] + w
                heapq.heappush(tas, (d[v], v))
    return d
\end{lstlisting}
\textbf{Complexité.} Chaque arête provoque au plus un ajout dans le tas, qui contient donc au
plus $E$ éléments~; chaque opération de tas coûte $O(\log E)=O(\log V)$ (car $E\le V^2$).
D'où $\boxed{O\big((V+E)\log V\big)}$.

\emph{Le test \code{if vus[u]: continue} est indispensable~: on n'enlève pas les entrées
périmées du tas, on les ignore lorsqu'on les extrait. C'est la version « paresseuse »,
la plus simple à écrire correctement.}

\Q \textbf{Invariant de Dijkstra.}
\begin{quote}\itshape
Au moment où un sommet $u$ est marqué (\code{vus[u] = True}), $\texttt{d[u]}$ est la distance
de $s$ à $u$~; de plus, pour tout sommet $v$ non encore marqué, $\texttt{d[v]}$ est le poids du
plus court chemin de $s$ à $v$ \textbf{dont tous les sommets intermédiaires sont marqués}.
\end{quote}
\emph{Démonstration de la première partie, par l'absurde.} Supposons que $u$ soit le premier
sommet marqué pour lequel $\texttt{d[u]}>\delta(s,u)$, où $\delta$ désigne la vraie distance.
Soit $\gamma$ un plus court chemin de $s$ à $u$, et soit $y$ le premier sommet de $\gamma$
non encore marqué au moment où l'on marque $u$ (il en existe, car $u$ l'est), et $x$ son
prédécesseur sur $\gamma$ ($x$ est marqué, et $x$ existe car $s$ est marqué en premier). Par
minimalité de $u$, $\texttt{d[x]}=\delta(s,x)$~; comme l'arête $(x,y)$ a été relâchée lors du
marquage de $x$, on a $\texttt{d[y]}\le\delta(s,x)+w(x,y)=\delta(s,y)$. Or $y$ est sur un plus
court chemin vers $u$ et \textbf{les poids sont positifs}, donc
$\delta(s,y)\le\delta(s,u)$. On obtient
\[ \texttt{d[y]}\ \le\ \delta(s,y)\ \le\ \delta(s,u)\ <\ \texttt{d[u]}, \]
ce qui contredit le fait que l'algorithme a extrait $u$ — de clé minimale — avant $y$.
\hfill$\square$

\textbf{Où sert la positivité~?} Exactement à l'inégalité $\delta(s,y)\le\delta(s,u)$~: un
préfixe d'un plus court chemin n'est plus court que le chemin entier que si les poids sont
$\ge 0$. Avec une arête de poids négatif, Dijkstra peut marquer définitivement un sommet à
une valeur trop grande et ne jamais se corriger.

\Q
\begin{lstlisting}
def bellman_ford(V, aretes, s):
    INF = float('inf')
    d = [INF] * V
    d[s] = 0
    for _ in range(V - 1):
        change = False
        for (u, v, w) in aretes:
            if d[u] + w < d[v]:
                d[v] = d[u] + w
                change = True
        if not change:
            break                              # arret anticipe
    for (u, v, w) in aretes:
        if d[u] + w < d[v]:
            return None                        # circuit de poids negatif
    return d
\end{lstlisting}
\textbf{Complexité.} $\Theta(V\!E)$ dans le pire des cas, $\Theta(E)$ en espace. C'est plus
coûteux que Dijkstra, mais l'algorithme tolère les poids négatifs et — surtout — il est
\emph{réparti}~: chaque sommet peut effectuer sa relaxation en ne connaissant que ses
voisins. C'est ce qui en fait le c\oe ur des protocoles à vecteur de distance.

\Q \textbf{Invariant.} \emph{Après $i$ tours de la boucle externe, pour tout sommet $v$,
$\texttt{d[v]}\le \delta_i(s,v)$, où $\delta_i(s,v)$ est le poids minimal d'un chemin de $s$ à
$v$ utilisant au plus $i$ arêtes ($+\infty$ s'il n'en existe pas).}

\emph{Initialisation.} Pour $i=0$~: $\delta_0(s,s)=0=\texttt{d[s]}$ et $\delta_0(s,v)=+\infty$
pour $v\ne s$.

\emph{Hérédité.} Soit $\gamma$ un chemin de $s$ à $v$ à au plus $i+1$ arêtes de poids
$\delta_{i+1}(s,v)$. Si $\gamma$ a au plus $i$ arêtes, la conclusion découle de l'hypothèse
de récurrence et du fait que $\texttt{d[v]}$ ne fait que décroître. Sinon, notons $u$
l'avant-dernier sommet de $\gamma$~: le préfixe de $\gamma$ jusqu'à $u$ a au plus $i$ arêtes
et est optimal parmi ceux-là, donc à la fin du tour $i$ on a $\texttt{d[u]}\le\delta_i(s,u)$.
Lors du tour $i+1$, l'arc $(u,v,w)$ est relâché, d'où
$\texttt{d[v]}\le\texttt{d[u]}+w\le\delta_i(s,u)+w=\delta_{i+1}(s,v)$. \hfill$\square$

\textbf{Correction et nombre de tours.} En l'absence de circuit de poids négatif, un plus
court chemin est \emph{élémentaire}, donc comporte au plus $V-1$ arêtes~: après $V-1$ tours,
$\texttt{d[v]}\le\delta_{V-1}(s,v)=\delta(s,v)$. L'inégalité inverse est immédiate puisque
toute valeur de $\texttt{d[v]}$ est le poids d'un chemin réel. D'où l'égalité. Enfin, si un
tour supplémentaire produit encore une amélioration, c'est qu'il existe un chemin à $V$
arêtes strictement meilleur que tout chemin à $V-1$ arêtes~: un tel chemin contient un
circuit, nécessairement de poids strictement négatif.

\Q \textbf{(Comptage à l'infini.)} Sur la ligne $A - B - C$, à l'équilibre $B$ annonce
$d(C)=1$ et $A$ annonce $d(C)=2$. Le lien $B-C$ tombe~: $B$ passe $d(C)$ à $+\infty$, mais il
reçoit alors de $A$ l'annonce « je sais atteindre $C$ en $2$ ». $B$ en déduit $d(C)=3$ (via
$A$) et l'annonce. $A$ recalcule alors $d(C)=4$ (via $B$), etc.

Le mécanisme est clair~: $A$ a appris sa route \emph{par $B$} mais l'annonce à $B$ sans
préciser cette provenance. Chaque routeur croit qu'un chemin subsiste chez l'autre~; les
distances croissent d'un pas à chaque échange, indéfiniment. C'est le \textbf{comptage à
l'infini}, défaut structurel des algorithmes à vecteur de distance — Bellman--Ford réparti
converge, mais lentement, et transite par des états incohérents.

\textbf{Contre-mesures.} RIP en emploie deux.
\begin{itemize}[leftmargin=1.4em,itemsep=1pt]
\item Le \emph{split horizon}~: un routeur n'annonce pas à un voisin une route apprise de ce
voisin (variante \emph{poison reverse}~: il l'annonce à distance infinie). Cela règle le cas
à deux routeurs, mais pas les boucles plus longues.
\item Un plafond~: la distance $16$ est déclarée infinie. Le comptage s'arrête donc en au
plus $16$ échanges — au prix d'une limitation drastique du diamètre du réseau, ce qui
cantonne RIP aux petits domaines.
\end{itemize}
C'est précisément pour éviter ces pathologies qu'OSPF a adopté l'approche à \emph{état de
liens}~: chaque routeur diffuse la topologie complète et exécute Dijkstra localement, au prix
d'un coût mémoire et d'un trafic de diffusion supérieurs.

\Q On applique Dijkstra depuis le sommet $0$. On note entre crochets le tableau
$\texttt{d}$ (avec $\infty$ pour l'infini) et l'on souligne le sommet extrait.

\begin{center}\small
\begin{tabular}{c|c|l}
Étape & Sommet extrait & Tableau $d=[d_0,\dots,d_5]$ après relaxation\\\hline
0 & --- & $[0,\infty,\infty,\infty,\infty,\infty]$\\
1 & $0$ & $[0,7,9,\infty,\infty,14]$\\
2 & $1$ ($d=7$) & $[0,7,9,22,\infty,14]$\\
3 & $2$ ($d=9$) & $[0,7,9,20,\infty,11]$\\
4 & $5$ ($d=11$) & $[0,7,9,20,20,11]$\\
5 & $3$ ($d=20$) & $[0,7,9,20,20,11]$\\
6 & $4$ ($d=20$) & $[0,7,9,20,20,11]$
\end{tabular}
\end{center}
D'où $\boxed{d=[0,7,9,20,20,11]}$.

Un plus court chemin de $0$ à $4$ se lit en remontant les prédécesseurs~: $4$ a été atteint
depuis $5$ ($11+9=20$), $5$ depuis $2$ ($9+2=11$), $2$ depuis $0$ ($9$). D'où
\[ \boxed{0 \to 2 \to 5 \to 4}\qquad\text{de poids } 9+2+9=20. \]
On notera que le chemin direct $0\to 5$ coûte $14$ alors que le détour $0\to2\to5$ ne coûte
que $11$~: l'arête la plus « courte » n'est pas toujours la bonne, ce qui justifie
l'algorithme. Bellman--Ford donne évidemment le même tableau.

%=======================================================================
\section*{Partie IV --- TCP : fiabilité, temporisateur et congestion}
%=======================================================================

\Q
\begin{lstlisting}
M = 1 << 32

def avant(a, b):
    return ((b - a) % M) < (1 << 31) and a != b
\end{lstlisting}
\textbf{Démonstration.} Supposons que tous les numéros en circulation appartiennent à
$\{c,c+1,\dots,c+W-1\}$ modulo $M$, avec $W<2^{31}$. Soient $a=c+i$ et $b=c+j$ deux d'entre
eux, $0\le i,j\le W-1$, $i\ne j$. Alors $(b-a)\bmod M = (j-i)\bmod M$ avec
$-(W-1)\le j-i\le W-1$.
\begin{itemize}[leftmargin=1.4em,itemsep=1pt]
\item Si $j>i$ (c'est-à-dire $a$ précède $b$), alors $0<j-i\le W-1<2^{31}$ et
$(b-a)\bmod M=j-i<2^{31}$~: la fonction renvoie \code{True}.
\item Si $j<i$, alors $(b-a)\bmod M=M+(j-i)\ge M-(W-1)>M-2^{31}=2^{31}$~: la fonction
renvoie \code{False}.
\end{itemize}
Le verdict est donc correct. \hfill$\square$

\textbf{Ambiguïté pour $W=2^{31}$.} Prenons $a=0$ et $b=2^{31}$. Alors
$(b-a)\bmod M=2^{31}\not<2^{31}$~: la fonction déclare que $a$ ne précède pas $b$. Mais
$(a-b)\bmod M=2^{31}$ également~: elle déclare aussi que $b$ ne précède pas $a$. Les deux
numéros sont à distance égale dans les deux sens sur le cercle~: la question « lequel vient
avant~? » n'a tout simplement pas de réponse. C'est pourquoi TCP impose une fenêtre
strictement inférieure à $2^{31}$ octets.

\Q
\begin{lstlisting}
def fusionner(segments):
    L = sorted(segments)                       # tri lexicographique : par debut
    res = []
    for (d, f) in L:
        if res and d <= res[-1][1]:            # chevauchement ou contiguite
            res[-1] = (res[-1][0], max(res[-1][1], f))
        else:
            res.append((d, f))
    return res
\end{lstlisting}
\textbf{Complexité.} $\Theta(k\log k)$ pour le tri, puis $\Theta(k)$ pour le balayage~: total
$\Theta(k\log k)$. C'est optimal dans le modèle des comparaisons, puisque le problème permet
de trier ($k$ intervalles disjoints en entrée ressortent triés).

\Q \textbf{Invariant.}
\begin{quote}\itshape
Avant le traitement de $L[i]$, la liste \code{res} contient les intervalles maximaux de la
réunion $\bigcup_{j<i}L[j]$, triés par début croissant, deux à deux disjoints et non
adjacents~; de plus, le début de $L[i]$ est supérieur ou égal au début du dernier élément de
\code{res}.
\end{quote}
\emph{Initialisation.} Pour $i=0$, \code{res} est vide~: trivial.

\emph{Hérédité.} Soit $(d,f)=L[i]$ et $(d_0,f_0)$ le dernier élément de \code{res}. Comme $L$
est trié par début croissant, $d\ge d_0$, et tout intervalle de \code{res} autre que le
dernier se termine avant $d_0\le d$~: seul le dernier peut rencontrer $(d,f)$. Deux cas.
\begin{itemize}[leftmargin=1.4em,itemsep=1pt]
\item Si $d\le f_0$, les deux intervalles se rencontrent ou sont contigus~; comme
$d_0\le d$, leur réunion est exactement $[\,d_0,\max(f_0,f)\,[$, ce qu'affecte l'algorithme.
Le nombre d'intervalles ne change pas, l'invariant est préservé.
\item Si $d>f_0$, l'intervalle $(d,f)$ est disjoint de tous les précédents \emph{et} non
adjacent au dernier~; on l'ajoute en queue, ce qui préserve le tri et la disjonction.
\end{itemize}
\hfill$\square$

\emph{Conclusion.} À la fin, \code{res} est la décomposition en intervalles maximaux de la
réunion de tous les segments. Cette décomposition est \textbf{unique}~: c'est la partition
en composantes connexes de la réunion, vue comme partie de $\mathbb{Z}$.

\emph{Vérification expérimentale~: sur $2\,000$ jeux de segments aléatoires, la réunion des
intervalles renvoyés a toujours coïncidé, ensemble d'entiers par ensemble d'entiers, avec la
réunion des segments fournis, et la sortie était toujours strictement croissante.}

\Q
\begin{lstlisting}
def rto_initial(R, G=0.001):
    SRTT = R
    RTTVAR = R / 2
    return SRTT, RTTVAR, max(1.0, SRTT + max(G, 4 * RTTVAR))

def rto_suivant(SRTT, RTTVAR, R, alpha=1/8, beta=1/4, G=0.001):
    RTTVAR = (1 - beta) * RTTVAR + beta * abs(SRTT - R)   # D'ABORD la variance
    SRTT   = (1 - alpha) * SRTT + alpha * R               # ENSUITE la moyenne
    return SRTT, RTTVAR, max(1.0, SRTT + max(G, 4 * RTTVAR))
\end{lstlisting}
\textbf{Pourquoi cet ordre~?} La mise à jour de $\mathrm{RTTVAR}$ utilise l'écart
$|\mathrm{SRTT}-R'|$ entre l'\emph{ancienne} estimation et la nouvelle mesure. Si l'on
mettait d'abord $\mathrm{SRTT}$ à jour, cet écart serait calculé avec une moyenne \emph{déjà
attirée} vers $R'$, donc systématiquement sous-évalué~: l'estimateur de dispersion serait
biaisé vers le bas, le $\mathrm{RTO}$ trop court, et les retransmissions inutiles se
multiplieraient. Le \rfc{6298} écrit d'ailleurs explicitement que l'ordre \emph{doit} être
respecté.

\Q \textbf{Démonstration par récurrence sur $n$.} La relation est
$S_n=(1-\alpha)S_{n-1}+\alpha R_n$ pour $n\ge 1$, avec $S_0=R_0$.

\emph{Initialisation ($n=0$).} La formule donne $(1-\alpha)^0R_0=R_0=S_0$, la somme étant
vide.

\emph{Hérédité.} Supposons la formule vraie au rang $n$. Alors
\[
S_{n+1}=(1-\alpha)S_n+\alpha R_{n+1}
=(1-\alpha)\Big[(1-\alpha)^nR_0+\alpha\sum_{i=1}^{n}(1-\alpha)^{n-i}R_i\Big]+\alpha R_{n+1}
\]
\[
=(1-\alpha)^{n+1}R_0+\alpha\sum_{i=1}^{n}(1-\alpha)^{n+1-i}R_i+\alpha(1-\alpha)^{0}R_{n+1}
=(1-\alpha)^{n+1}R_0+\alpha\sum_{i=1}^{n+1}(1-\alpha)^{n+1-i}R_i .
\]
\hfill$\square$

\textbf{Somme des coefficients.} En posant $j=n-i$, la somme géométrique donne
\[ (1-\alpha)^n+\alpha\sum_{i=1}^{n}(1-\alpha)^{n-i}
 = (1-\alpha)^n+\alpha\sum_{j=0}^{n-1}(1-\alpha)^{j}
 = (1-\alpha)^n+\alpha\cdot\frac{1-(1-\alpha)^n}{\alpha}
 = 1 . \]
$S_n$ est donc bien une \textbf{moyenne pondérée} des mesures, les poids décroissant
géométriquement vers le passé~: c'est une moyenne mobile exponentielle. Cela garantit en
particulier $\min_i R_i\le S_n\le\max_i R_i$.

\emph{Vérification~: sur $500$ suites aléatoires, la forme close et l'itération coïncident à
$10^{-12}$ près.}

\Q Les poids de la question 33, lus du plus récent au plus ancien, sont
$\alpha,\ \alpha(1-\alpha),\ \alpha(1-\alpha)^2,\dots$ On calcule, pour $|1-\alpha|<1$,
\[ \mu=\sum_{j\ge 0}(j+1)\,\alpha(1-\alpha)^{j}
     =\alpha\sum_{j\ge 0}(j+1)x^{j}\Big|_{x=1-\alpha}
     =\frac{\alpha}{(1-x)^2}\Big|_{x=1-\alpha}
     =\frac{\alpha}{\alpha^2}=\boxed{\frac{1}{\alpha}} \]
(en dérivant terme à terme la série géométrique $\sum x^{j}=\frac{1}{1-x}$, licite sur
$]-1,1[$). Pour $\alpha=1/8$~: $\mu=\boxed{8}$ échantillons.

Pour le second point, on cherche le plus petit $j$ tel que
$\alpha(1-\alpha)^{j}<\tfrac{1}{10}\alpha$, soit $(7/8)^{j}<1/10$, c'est-à-dire
$j>\ln 10/\ln(8/7)\approx 2{,}3026/0{,}1335\approx 17{,}24$. D'où $\boxed{j=18}$.

\textbf{Interprétation.} L'estimateur « se souvient » d'environ $8$ mesures et a
pratiquement oublié tout ce qui date de plus de $18$ aller-retours. C'est un compromis
assumé~: assez de mémoire pour lisser le bruit, assez peu pour suivre un changement de route
ou de charge en quelques secondes.

\Q \textbf{Algorithme de Karn.} Un segment retransmis donne lieu à \emph{deux} émissions et
à un seul accusé de réception, sans que l'on sache à laquelle il répond. Si l'on mesure $R$
depuis la \emph{seconde} émission alors que l'accusé répondait à la première, on obtient une
valeur beaucoup trop \textbf{petite}. L'estimateur diminue, le $\mathrm{RTO}$ diminue, donc
davantage de segments expirent et sont retransmis — d'où de nouvelles sous-estimations~: le
biais s'auto-entretient et peut provoquer un effondrement par congestion. Karn et Partridge
(1987) imposent donc de ne jamais échantillonner sur un segment retransmis, sauf si
l'estampille temporelle lève l'ambiguïté.

\textbf{Repli exponentiel.} Après $n$ expirations, les attentes successives valent
$T_0,2T_0,\dots,2^{n-1}T_0$, d'où le temps cumulé
\[ \sum_{i=0}^{n-1}2^{i}T_0 \;=\; T_0\,(2^{n}-1). \]
Pour $T_0=1$~s et $n=5$~: $\boxed{31\ \text{secondes}}$. Ce doublement est la réponse
canonique à une congestion~: en cas de doute, l'émetteur se retire exponentiellement du
réseau plutôt que d'insister.

\Q Les transmissions étant indépendantes et chacune réussissant avec probabilité $1-p$,
$N$ est le rang du premier succès dans une suite d'épreuves de Bernoulli~: $N$ suit une
\textbf{loi géométrique} de paramètre $1-p$, à valeurs dans $\mathbb{N}^{*}$~:
\[ \mathbb{P}(N=k)=p^{\,k-1}(1-p),\qquad \mathbb{P}(N>k)=p^{\,k},\qquad
   \mathbb{E}[N]=\frac{1}{1-p}. \]
($\mathbb{P}(N>k)$ s'obtient directement~: les $k$ premières transmissions échouent, ce qui a
probabilité $p^{k}$ par indépendance.)

\textbf{Terminaison presque sûre.} L'événement « ne jamais aboutir » est l'intersection
décroissante des $\{N>k\}$, donc par continuité décroissante
\[ \mathbb{P}(N=+\infty)=\lim_{k\to+\infty}p^{k}=0 \qquad\text{car } 0<p<1. \]
\hfill$\square$

Pour $p=0{,}1$~: $\mathbb{E}[N]=1/0{,}9\approx\boxed{1{,}11}$ transmission, et
$\mathbb{P}(N>3)=10^{-3}$. Le pire des cas n'est en revanche pas borné~: c'est un algorithme
de type Las Vegas, dont on ne peut analyser que le comportement moyen.

\Q Pendant un cycle, la fenêtre prend successivement les valeurs
$\frac{W_{\max}}{2},\ \frac{W_{\max}}{2}+1,\ \dots,\ W_{\max}-1$, à raison d'un
aller-retour par valeur~; le cycle dure donc $W_{\max}/2$ aller-retours. En posant
$W=W_{\max}$, le nombre de segments émis est
\[ N_{\text{cycle}}=\sum_{i=W/2}^{W-1} i
 =\sum_{i=0}^{W-1}i-\sum_{i=0}^{W/2-1}i
 =\frac{(W-1)W}{2}-\frac{(W/2-1)(W/2)}{2}, \]
soit, en développant,
\[ N_{\text{cycle}}=\frac{W^2-W}{2}-\frac{W^2/4-W/2}{2}
 =\boxed{\frac{3W^{2}}{8}-\frac{W}{4}} . \]
La fenêtre moyenne est le quotient par le nombre d'aller-retours~:
\[ \overline{W}=\frac{N_{\text{cycle}}}{W/2}
 =\frac{2}{W}\Big(\frac{3W^2}{8}-\frac{W}{4}\Big)=\frac{3W}{4}-\frac{1}{2}
 =\boxed{\frac{3W-2}{4}} . \]
\hfill$\square$

\emph{Vérification~: pour $W=8,\,16,\,64,\,100$ on obtient respectivement
$N_{\text{cycle}}=22,\ 92,\ 1520,\ 3725$ et $\overline{W}=5{,}5\ ;\ 11{,}5\ ;\ 47{,}5\ ;\
74{,}5$, en accord exact avec les formules.}

\Q Une perte par cycle et une probabilité de perte $p$ par segment donnent
$p=1/N_{\text{cycle}}$, soit
\[ \frac{1}{p}=\frac{3W^{2}}{8}-\frac{W}{4}\ \sim\ \frac{3W^{2}}{8}
   \qquad\text{quand } W\to+\infty, \]
d'où $W\approx\sqrt{\dfrac{8}{3p}}$. Le débit est le nombre d'octets émis par unité de temps,
c'est-à-dire $\overline{W}\cdot\mathrm{MSS}/\mathrm{RTT}$ avec
$\overline{W}\sim\frac{3}{4}W$~:
\[ \text{débit}\ \approx\ \frac{3}{4}\sqrt{\frac{8}{3p}}\cdot
   \frac{\mathrm{MSS}}{\mathrm{RTT}}
  =\frac{\mathrm{MSS}}{\mathrm{RTT}}\sqrt{\frac{9}{16}\cdot\frac{8}{3p}}
  =\frac{\mathrm{MSS}}{\mathrm{RTT}}\sqrt{\frac{3}{2p}}
  \approx\frac{1{,}2247\ \mathrm{MSS}}{\mathrm{RTT}\sqrt{p}} . \]
\hfill$\square$

C'est la \textbf{formule de Mathis} (1997), l'un des résultats les plus cités de la
littérature réseau~: elle donne, à partir de trois paramètres seulement, le débit d'une
connexion TCP en régime établi.

\Q Avec $\mathrm{MSS}=1460$ octets et $\mathrm{RTT}=0{,}100$~s~:
\begin{center}
\begin{tabular}{c|c|c|c}
$p$ & $W_{\max}\approx\sqrt{8/(3p)}$ & débit (octets/s) & débit (Mbit/s)\\\hline
$10^{-2}$ & $16{,}3$ & $1{,}79\times 10^{5}$ & $1{,}43$\\
$10^{-4}$ & $163{,}3$ & $1{,}79\times 10^{6}$ & $14{,}3$
\end{tabular}
\end{center}

\Q \textbf{(a)} Le débit varie comme $p^{-1/2}$~: diviser le taux de perte par $4$ ne
multiplie le débit que par $2$. Les rendements sont donc \emph{décroissants}~; améliorer la
qualité d'un lien déjà bon rapporte peu. C'est un argument classique en faveur des
mécanismes de retransmission locale sur les liens sans fil.

\textbf{(b)} Le débit varie comme $1/\mathrm{RTT}$~: à taux de perte égal, une connexion de
$\mathrm{RTT}$ deux fois plus court obtient un débit deux fois plus grand. Deux connexions
partageant le même goulet d'étranglement ne se partagent donc \textbf{pas} équitablement la
bande passante~: la plus proche écrase la plus lointaine. C'est le fameux biais de TCP contre
les connexions à long délai, l'une des principales motivations des algorithmes de congestion
modernes (CUBIC, dont la croissance ne dépend plus du $\mathrm{RTT}$~; BBR, qui estime
directement le débit et le délai plutôt que d'inférer la congestion des pertes).

%=======================================================================
\section*{Partie V --- HTTP, cache et HTTPS}
%=======================================================================

\Q
\begin{lstlisting}
def analyser_requete(texte):
    i = texte.find('\r\n\r\n')
    if i < 0:
        return None                          # en-tete incomplet
    lignes = texte[:i].split('\r\n')
    methode, cible, version = lignes[0].split(' ')
    entetes = {}
    for ligne in lignes[1:]:
        j = ligne.index(':')
        entetes[ligne[:j].lower()] = ligne[j+1:].strip()
    return (methode, cible, version, entetes, texte[i+4:])
\end{lstlisting}
Complexité $\Theta(n)$ où $n$ est la longueur du texte~: la recherche, la découpe et le
parcours des lignes sont linéaires. La mise en minuscules des noms est imposée par le
\rfc{9110}~: les noms d'en-têtes sont insensibles à la casse.

\emph{Sur l'exemple du sujet, on obtient}
\begin{lstlisting}
('GET', '/index.html?p=1', 'HTTP/1.1',
 {'host': 'cpgeacademy.org', 'content-length': '5'}, 'HELLO')
\end{lstlisting}

\Q On cherche la première occurrence du mot $u=\texttt{CRLFCRLF}$ de longueur $4$. L'automate
a pour états $\{0,1,2,3\}$, l'état $q$ signifiant « les $q$ derniers caractères lus forment
le préfixe de longueur $q$ de $u$ » ($4$ est l'état d'acceptation, où l'on s'arrête).

\begin{center}\small
\begin{tabular}{c|c|c|c|l}
état & \texttt{'\textbackslash r'} & \texttt{'\textbackslash n'} & autre &
 signification\\\hline
$0$ & $1$ & $0$ & $0$ & rien de reconnu\\
$1$ & $1$ & $2$ & $0$ & \texttt{CR} lu\\
$2$ & $3$ & $0$ & $0$ & \texttt{CRLF} lu\\
$3$ & $1$ & \textbf{accepter} & $0$ & \texttt{CRLFCR} lu
\end{tabular}
\end{center}
La seule subtilité est la transition $1\xrightarrow{\texttt{CR}}1$ et
$3\xrightarrow{\texttt{CR}}1$~: après un \texttt{CR} inattendu, ce \texttt{CR} peut lui-même
amorcer une nouvelle occurrence, il ne faut donc pas retomber en $0$.
\begin{lstlisting}
def fin_entete(flux):
    etat = 0
    for k, c in enumerate(flux):
        if   etat == 0: etat = 1 if c == '\r' else 0
        elif etat == 1: etat = 2 if c == '\n' else (1 if c == '\r' else 0)
        elif etat == 2: etat = 3 if c == '\r' else 0
        elif etat == 3:
            if c == '\n':
                return k + 1
            etat = 1 if c == '\r' else 0
    return -1
\end{lstlisting}

\Q \textbf{Invariant.}
\begin{quote}\itshape
Après la lecture des $k$ premiers caractères sans acceptation, \code{etat} est la longueur du
plus long suffixe du préfixe lu qui soit un préfixe strict de $u=\texttt{CRLFCRLF}$.
\end{quote}
\emph{Initialisation.} Avant toute lecture, le seul suffixe est le mot vide~: $\texttt{etat}=0$.

\emph{Hérédité.} Soit $q$ l'état courant et $c$ le caractère lu. Le plus long suffixe du
nouveau texte qui soit préfixe de $u$ est nécessairement de longueur $\le q+1$. S'il est de
longueur $q+1$, c'est que $u_{q}=c$~: c'est exactement la transition « avancer ». Sinon, il
faut prendre le plus long \emph{bord} de $u_{0..q-1}$ suivi de $c$. Ici $u=\texttt{CRLFCRLF}$
et $u_{0..2}=\texttt{CRLF}$ n'a pour bords que le mot vide, sauf $u_{0..0}=\texttt{CR}$ dont
le seul bord est vide~: en pratique, si $c=\texttt{CR}$ on retombe en $1$, sinon en $0$. Ce
sont précisément les transitions écrites. \hfill$\square$

\emph{Conséquences.} Chaque caractère est lu \textbf{une seule fois} et traité en $O(1)$~:
l'analyse est en $\Theta(n)$, sans retour en arrière ni mémoire tampon. C'est essentiel pour
un serveur qui traite un flux au fil de l'eau, sans pouvoir revenir sur les octets déjà
consommés.

\textbf{Vérification sur} \code{"\\r\\n\\r\\r\\n\\r\\n"} (longueur $7$)~:
\begin{center}\small
\begin{tabular}{c|cccccccc}
$k$ & \multicolumn{1}{c}{début} & $0$ & $1$ & $2$ & $3$ & $4$ & $5$ & $6$\\
caractère & --- & \texttt{CR} & \texttt{LF} & \texttt{CR} & \texttt{CR} & \texttt{LF} &
 \texttt{CR} & \texttt{LF}\\\hline
état après & $0$ & $1$ & $2$ & $3$ & $1$ & $2$ & $3$ & \textbf{accept.}
\end{tabular}
\end{center}
La fonction renvoie $7$. On vérifie l'invariant au rang $k=3$~: le texte lu est
\texttt{CRLFCR}, dont le plus long suffixe préfixe de $u$ après lecture du second \texttt{CR}
est \texttt{CR}, de longueur $1$ — c'est bien l'état obtenu. Une implantation naïve qui
serait retombée en $0$ aurait manqué l'occurrence.

\Q
\begin{lstlisting}
def decoder_chunked(s):
    res = []
    i = 0
    while True:
        j = s.index('\r\n', i)               # fin de la ligne de taille
        n = int(s[i:j], 16)                  # taille en hexadecimal
        if n == 0:
            return ''.join(res)
        res.append(s[j+2 : j+2+n])
        i = j + 2 + n + 2                    # on saute le bloc et son CRLF
\end{lstlisting}
\textbf{Variant.} Posons $V=\texttt{len(s)} - i$. À chaque tour, $i$ augmente d'au moins
$\texttt{len(hex)}+2+n+2\ge 5>0$~: $V$ est un entier strictement décroissant, minoré par $0$
tant que la chaîne est bien formée. La boucle se termine donc, au plus tard en épuisant la
chaîne (ce qui lèverait alors une exception, comportement attendu sur une entrée malformée).

\textbf{Complexité.} $\Theta(n)$~: chaque caractère est examiné au plus deux fois (une par
\code{index}, une par la copie), et la concaténation finale par \code{join} est linéaire.
On évitera soigneusement \code{res = res + bloc} dans la boucle, qui donnerait $\Theta(n^2)$.

\emph{Exemple~:} le corps \code{"4\\r\\nWiki\\r\\n5\\r\\npedia\\r\\nE\\r\\n in\\r\\n\\r\\nchunks.\\r\\n0\\r\\n\\r\\n"}
se décode en \code{"Wikipedia in\\r\\n\\r\\nchunks."} — noter que le troisième bloc, de
taille $\texttt{E}_{16}=14$, contient lui-même des \texttt{CRLF}~: c'est tout l'intérêt du
préfixe de longueur.

\Q La difficulté est d'obtenir $O(1)$ \emph{à la fois} pour la recherche et pour la mise à
jour de l'ordre d'ancienneté. On combine donc~:
\begin{itemize}[leftmargin=1.4em,itemsep=1pt]
\item un \textbf{dictionnaire} clé $\mapsto$ n\oe ud, pour la recherche en $O(1)$ moyen~;
\item une \textbf{liste doublement chaînée} des clés, de la plus récemment utilisée (en tête)
à la plus ancienne (en queue), avec deux n\oe uds sentinelles pour éviter les cas
particuliers.
\end{itemize}
Un n\oe ud est une liste \code{[cle, precedent, suivant]}.
\begin{lstlisting}
class LRU:
    def __init__(self, cap):
        self.cap = cap
        self.d = {}
        self.tete  = ['GARDE', None, None]         # sentinelles
        self.queue = ['GARDE', self.tete, None]
        self.tete[2] = self.queue

    def _retirer(self, x):
        x[1][2] = x[2]
        x[2][1] = x[1]

    def _en_tete(self, x):
        x[1] = self.tete
        x[2] = self.tete[2]
        self.tete[2][1] = x
        self.tete[2] = x

    def acceder(self, cle):
        if cle in self.d:                          # succes
            x = self.d[cle]
            self._retirer(x)
            self._en_tete(x)
            return True
        if len(self.d) == self.cap:                # eviction du plus ancien
            v = self.queue[1]
            self._retirer(v)
            del self.d[v[0]]
        x = [cle, None, None]
        self._en_tete(x)
        self.d[cle] = x
        return False
\end{lstlisting}
\textbf{Invariant.}
\begin{quote}\itshape
À tout instant, la liste chaînée contient exactement les clés du dictionnaire, dans l'ordre
décroissant de la date de leur dernier accès~; le dictionnaire associe à chaque clé le n\oe ud
qui la porte~; et le nombre de clés n'excède jamais \code{cap}.
\end{quote}
Chaque opération de \code{acceder} — test d'appartenance, retrait, insertion en tête,
suppression dans le dictionnaire — se fait en $O(1)$ moyen, le retrait en $O(1)$ exact grâce
au \emph{double} chaînage (sans lui, retirer un n\oe ud demanderait de retrouver son
prédécesseur, en $O(C)$). D'où un accès en $\boxed{O(1)}$ en moyenne, et un espace $O(C)$.

\emph{Vérification expérimentale~: sur $3\,000$ suites d'accès aléatoires et des capacités de
$1$ à $5$, cette classe et une implantation naïve de référence (liste Python réordonnée à
chaque accès) ont toujours donné le même nombre de succès.}

\Q
\begin{center}\small
\begin{tabular}{c|c|l|c}
Accès & Résultat & Cache après (du plus récent au plus ancien) & Éviction\\\hline
$a$ & échec & $[a]$ & ---\\
$b$ & échec & $[b,a]$ & ---\\
$c$ & échec & $[c,b,a]$ & ---\\
$a$ & \textbf{succès} & $[a,c,b]$ & ---\\
$d$ & échec & $[d,a,c]$ & $b$\\
$b$ & échec & $[b,d,a]$ & $c$\\
$a$ & \textbf{succès} & $[a,b,d]$ & ---\\
$c$ & échec & $[c,a,b]$ & $d$\\
$d$ & échec & $[d,c,a]$ & $b$\\
$a$ & \textbf{succès} & $[a,d,c]$ & ---
\end{tabular}
\end{center}
On compte $\boxed{3\ \text{succès sur }10}$, soit un taux de succès de $30\,\%$.

\emph{Commentaire.} L'objet $a$, régulièrement redemandé, est toujours conservé~: c'est le
comportement attendu. En revanche $b$, $c$ et $d$ s'évincent mutuellement en boucle. La
politique optimale hors ligne (algorithme de Belady, qui évince l'objet dont la prochaine
demande est la plus lointaine) ferait mieux~; mais elle suppose de connaître l'avenir, ce qui
est exclu en ligne.

\Q \textbf{Initialisation.} $m_0(j)=0$ pour tout $j$~: sans objet, on ne sert rien.

\textbf{Récurrence.} Un sous-ensemble optimal de $\{0,\dots,i\}$ de taille $\le j$ contient
ou non l'objet $i$~:
\[ m_{i+1}(j) \;=\;
\begin{cases}
m_i(j) & \text{si } j<t_i,\\[3pt]
\max\big(\,m_i(j),\ m_i(j-t_i)+f_i\,\big) & \text{si } j\ge t_i .
\end{cases} \]
Dans le second cas, le premier terme correspond au choix « ne pas prendre $i$ », le second à
« prendre $i$ », ce qui laisse une capacité $j-t_i$ pour les objets $0,\dots,i-1$. C'est bien
une partition des solutions, d'où le maximum.

\Q
\begin{lstlisting}
def cache_optimal(t, f, C):
    m = [0] * (C + 1)
    for i in range(len(t)):
        for j in range(C, t[i] - 1, -1):        # parcours DECROISSANT
            if m[j - t[i]] + f[i] > m[j]:
                m[j] = m[j - t[i]] + f[i]
    return m[C]
\end{lstlisting}
\textbf{Sens de parcours.} La case $j$ est calculée à partir de la case $j-t_i<j$. En
parcourant $j$ de $C$ vers $t_i$, la case $j-t_i$ n'a pas encore été modifiée au cours de
cette itération~: elle contient donc $m_i(j-t_i)$, valeur qui \emph{n'utilise pas} l'objet
$i$. L'objet est ainsi pris au plus une fois.

\textbf{Contre-exemple en parcours croissant.} Prenons un seul objet, $t=[3]$, $f=[10]$, et
$C=9$. En croissant~: $j=3$ donne $m[3]=10$~; puis $j=6$ lit $m[3]=10$ et donne $m[6]=20$~;
puis $j=9$ lit $m[6]$ et donne $m[9]=30$. L'algorithme renverrait $30$ en plaçant
\emph{trois fois} le même objet dans le cache, alors que la réponse correcte est $10$. (Le
parcours croissant résout en fait le sac à dos \emph{avec répétition}, un autre problème.)

\Q \textbf{Invariant.}
\begin{quote}\itshape
À la fin de la $i$-ième itération de la boucle externe, pour tout $j\in[\![0,C]\!]$,
$\texttt{m[j]}=m_i(j)$, c'est-à-dire le nombre maximal de requêtes servies par un sous-ensemble
de $\{0,\dots,i-1\}$ de taille totale au plus $j$.
\end{quote}
\emph{Initialisation.} Avant la première itération, $\texttt{m}$ est nul, c'est $m_0$.

\emph{Hérédité.} Supposons $\texttt{m}=m_i$ au début de l'itération traitant l'objet $i$.
Notons $\texttt{m}^{\text{fin}}$ le tableau à la fin de cette itération.

$(\le)$ Toute valeur écrite dans $\texttt{m}^{\text{fin}}[j]$ vaut $m_i(j-t_i)+f_i$, donc est
réalisée par un sous-ensemble de $\{0,\dots,i-1\}$ de taille $\le j-t_i$ auquel on ajoute
$i$~: c'est un sous-ensemble de $\{0,\dots,i\}$ de taille $\le j$. Les valeurs conservées
sont les $m_i(j)$, également réalisables. Donc
$\texttt{m}^{\text{fin}}[j]\le m_{i+1}(j)$.

$(\ge)$ Soit $S\subseteq\{0,\dots,i\}$ de taille $\le j$ réalisant $m_{i+1}(j)$. Si
$i\notin S$ alors $\sum_{k\in S}f_k\le m_i(j)\le\texttt{m}^{\text{fin}}[j]$, puisque le tableau
n'est jamais diminué. Si $i\in S$, alors $S\setminus\{i\}$ est de taille $\le j-t_i$, donc
$\sum_{k\in S}f_k\le m_i(j-t_i)+f_i$~; or, grâce au parcours décroissant, la case $j-t_i$
contenait bien $m_i(j-t_i)$ au moment de la lecture, et la mise à jour a porté la case $j$ à
au moins cette valeur. Donc $\texttt{m}^{\text{fin}}[j]\ge m_{i+1}(j)$.

D'où l'égalité, et l'invariant. \hfill$\square$

\emph{Terminaison.} Les deux boucles sont des boucles \code{for} sur des intervalles finis~:
la terminaison est immédiate, le nombre total d'itérations étant au plus $n(C+1)$.

\emph{Conclusion.} À la fin, $\texttt{m[C]}=m_n(C)=\mu(t,f,C)$.

\Q \textbf{Complexité.} Temps $\Theta(nC)$~; espace $\Theta(C)$ auxiliaire, contre
$\Theta(nC)$ si l'on conservait tous les $m_i$.

\textbf{Appartenance à NP.} Un certificat de l'instance « existe-t-il $S$ avec
$\sum_{i\in S}t_i\le C$ et $\sum_{i\in S}f_i\ge F$~? » est la liste des indices de $S$, de
taille $O(n\log n)$ bits donc polynomiale. Le vérificateur calcule les deux sommes ($2n$
additions d'entiers de taille au plus celle de l'entrée) et effectue deux comparaisons~: il
est polynomial. Donc le problème est dans \textbf{NP}. \hfill$\square$

\textbf{Pourquoi il n'y a pas contradiction.} La taille de l'\emph{entrée} n'est pas $C$ mais
le nombre de bits nécessaires pour l'écrire, soit $\Theta(\log C)$. La complexité
$\Theta(nC)=\Theta\big(n\,2^{\log_2 C}\big)$ est donc \emph{exponentielle} en la taille de
l'entrée. On dit l'algorithme \textbf{pseudo-polynomial}~: polynomial en la \emph{valeur} des
données, non en leur \emph{taille}.

\textbf{Tailles en octets, $C=10$~Go.} On aurait $C\approx 10^{10}$~: le tableau à lui seul
occuperait des dizaines de gigaoctets et le calcul demanderait $n\cdot 10^{10}$ opérations —
irréalisable. On change donc d'unité (par exemple le mégaoctet, ou le PGCD des tailles), ce
qui divise $C$ par autant, ou l'on renonce à l'optimum exact au profit de l'heuristique
gloutonne consistant à trier les objets par densité $f_i/t_i$ décroissante — laquelle donne,
en version fractionnaire, une $2$-approximation.

\Q Avec $t=[3,4,5,9,4]$, $f=[60,50,70,90,30]$, $C=10$~:
\[ \boxed{\mu=130},\qquad\text{atteint par } S^{\star}=\{0,2\}\ \ (\text{tailles } 3+5=8\le
10,\ \text{fréquences } 60+70=130). \]
C'est le \emph{seul} sous-ensemble optimal. On remarquera que l'objet $3$, de fréquence
maximale $90$, n'est pas retenu~: sa taille $9$ mobiliserait presque tout le cache pour un
gain moindre. C'est exactement ce que le glouton par fréquence décroissante manquerait
($90$ puis plus rien, soit $90<130$)~; le glouton par densité $f_i/t_i$, lui, donne
$70/5=14$, $60/3=20$, $50/4=12{,}5$, $30/4=7{,}5$, $90/9=10$, soit l'ordre $0,2,\dots$ et
retrouve ici l'optimum — mais ce n'est pas garanti en général.

\textbf{Déroulé pour $t=[3,4,5]$, $f=[60,50,70]$, $C=8$}~:
\[
\begin{array}{l|ccccccccc}
j & 0&1&2&3&4&5&6&7&8\\\hline
\text{initialisation} & 0&0&0&0&0&0&0&0&0\\
\text{après } (t_0,f_0)=(3,60) & 0&0&0&60&60&60&60&60&60\\
\text{après } (t_1,f_1)=(4,50) & 0&0&0&60&60&60&60&110&110\\
\text{après } (t_2,f_2)=(5,70) & 0&0&0&60&60&70&70&110&\mathbf{130}
\end{array}
\]
D'où $\mu([3,4,5],[60,50,70],8)=\boxed{130}$, réalisé par $\{0,2\}$ (taille $8$).

\subsection*{Requêtes SQL}

\Q
\begin{lstlisting}[style=sql]
SELECT methode, COUNT(*) AS nb, SUM(octets) AS total_octets
FROM journal
GROUP BY methode
ORDER BY nb DESC;
\end{lstlisting}

\Q
\begin{lstlisting}[style=sql]
SELECT chemin, COUNT(*) AS nb
FROM journal
GROUP BY chemin
ORDER BY nb DESC, chemin ASC
LIMIT 3;
\end{lstlisting}
Le second critère de tri (\code{chemin ASC}) rend le résultat déterministe en cas d'égalité~:
sans lui, deux exécutions pourraient renvoyer des lignes différentes.

\Q
\begin{lstlisting}[style=sql]
SELECT SUBSTR(horodatage, 12, 2) AS heure,
       COUNT(*) AS total,
       SUM(CASE WHEN statut BETWEEN 500 AND 599 THEN 1 ELSE 0 END) AS erreurs,
       ROUND(100.0 * SUM(CASE WHEN statut BETWEEN 500 AND 599 THEN 1 ELSE 0 END)
             / COUNT(*), 1) AS taux
FROM journal
GROUP BY heure
ORDER BY heure;
\end{lstlisting}
Deux points méritent attention~: la position $12$ isole les deux caractères de l'heure dans
\code{'AAAA-MM-JJ HH:MM:SS'} (les positions sont comptées à partir de $1$)~; et le facteur
\code{100.0} force une division \emph{flottante} — avec \code{100}, la division entière
renverrait $0$ pour tout taux inférieur à $100\,\%$.

\Q
\begin{lstlisting}[style=sql]
SELECT ip, COUNT(*) AS nb4xx
FROM journal
WHERE statut BETWEEN 400 AND 499
GROUP BY ip
HAVING COUNT(*) >= 3
ORDER BY nb4xx DESC;
\end{lstlisting}
On notera la différence de rôle entre \code{WHERE}, qui filtre les \emph{lignes} avant le
regroupement, et \code{HAVING}, qui filtre les \emph{groupes} après.

\Q Volume par pays~:
\begin{lstlisting}[style=sql]
SELECT c.pays, COUNT(*) AS nb, SUM(j.octets) AS octets
FROM journal j
JOIN client c ON j.ip = c.ip
GROUP BY c.pays
ORDER BY octets DESC;
\end{lstlisting}
Chemins jamais servis avec le statut $200$~:
\begin{lstlisting}[style=sql]
SELECT DISTINCT chemin
FROM journal j1
WHERE NOT EXISTS (SELECT 1 FROM journal j2
                  WHERE j2.chemin = j1.chemin AND j2.statut = 200)
ORDER BY chemin;
\end{lstlisting}
La sous-requête corrélée exprime directement la quantification universelle « pour toute ligne
de ce chemin, le statut n'est pas $200$ ». Une écriture équivalente et souvent plus efficace
utilise l'agrégation~:
\begin{lstlisting}[style=sql]
SELECT chemin FROM journal
GROUP BY chemin
HAVING SUM(CASE WHEN statut = 200 THEN 1 ELSE 0 END) = 0;
\end{lstlisting}

\subsection*{HTTPS}

\Q
\begin{lstlisting}
def puissance_mod(a, e, n):
    r = 1
    a = a % n
    while e > 0:
        if e % 2 == 1:
            r = (r * a) % n
        a = (a * a) % n
        e = e // 2
    return r
\end{lstlisting}
\textbf{Terminaison.} Le variant est $e$~: c'est un entier naturel, et tant que la boucle
continue on a $e\ge 1$, donc $\lfloor e/2\rfloor<e$. Il décroît strictement~: la boucle
s'arrête. \hfill$\square$

\textbf{Invariant.} Notons $a_0$ et $e_0$ les valeurs initiales.
\begin{quote}\itshape
Au début de chaque tour, $r\cdot a^{\,e}\equiv a_0^{\,e_0} \pmod n$.
\end{quote}
\emph{Initialisation.} $r=1$, $a\equiv a_0$, $e=e_0$~: l'égalité est immédiate.

\emph{Hérédité.} Notons $r',a',e'$ les valeurs après un tour. Si $e$ est pair, $r'=r$,
$a'\equiv a^2$, $e'=e/2$, d'où $r'a'^{e'}\equiv r\,(a^2)^{e/2}=r\,a^{e}$. Si $e$ est impair,
$r'\equiv ra$, $a'\equiv a^{2}$, $e'=(e-1)/2$, d'où
$r'a'^{e'}\equiv ra\cdot a^{e-1}=r\,a^{e}$. Dans les deux cas la quantité est conservée.

\emph{Conclusion.} À la sortie $e=0$, donc $r\equiv r\cdot a^{0}\equiv a_0^{\,e_0}$.
\hfill$\square$

\textbf{Nombre de multiplications.} Le nombre de tours est le nombre de bits de $e$, soit
$\lfloor\log_2 e\rfloor+1$. Chaque tour effectue un carré et, au plus, une multiplication~:
d'où au plus $2(\lfloor\log_2 e\rfloor+1)=2\lfloor\log_2 e\rfloor+2$ multiplications
modulaires. Le calcul naïf en demanderait $e-1$~: pour $e$ de $2048$ bits, on passe de
$2^{2048}$ à moins de $4100$ multiplications.

\emph{Vérification~: $300$ tirages aléatoires avec $n=2^{61}-1$ et des exposants jusqu'à
$2^{40}$ redonnent exactement \code{pow(a, e, n)}.}

\Q Le client calcule $K_C=B^{x}\bmod p$ et le serveur $K_S=A^{y}\bmod p$.
\begin{lstlisting}
K_client  = puissance_mod(B, x, p)
K_serveur = puissance_mod(A, y, p)
\end{lstlisting}
\textbf{Démonstration.} Par associativité de l'exponentiation dans le groupe
$(\mathbb{Z}/p\mathbb{Z})^{*}$~:
\[ K_C \equiv B^{x} \equiv (g^{y})^{x} \equiv g^{xy} \equiv (g^{x})^{y} \equiv A^{y}
\equiv K_S \pmod p . \]
\hfill$\square$ Le secret $g^{xy}$ n'a jamais transité sur le réseau, bien qu'aucun des deux
participants n'ait eu besoin de connaître l'exposant de l'autre.

\textbf{Application numérique} ($p=23$, $g=5$, $x=6$, $y=15$)~:
\[ A=5^{6}\bmod 23 = 15\,625 \bmod 23 = \boxed{8},\qquad
   B=5^{15}\bmod 23 = \boxed{19}, \]
\[ K = B^{x}\bmod 23 = 19^{6}\bmod 23 = \boxed{2}
    \qquad\text{et}\qquad A^{y}\bmod 23 = 8^{15}\bmod 23 = 2 . \]
(Détail pour $A$~: $5^2=25\equiv 2$, $5^4\equiv 4$, $5^6=5^4\cdot 5^2\equiv 4\cdot 2=8$.)

\Q \textbf{Algorithme naïf.} On énumère $k=0,1,2,\dots$ en calculant $g^{k}$ de proche en
proche ($g^{k+1}=g\cdot g^{k}$) jusqu'à trouver $A$. Comme l'ordre de $g$ divise $p-1$, on
s'arrête au plus tard après $p-1$ essais~: complexité $O(p)$ en temps, $O(1)$ en espace.

\textbf{Pas de bébé, pas de géant.} Posons $m=\lceil\sqrt{p-1}\,\rceil$ et écrivons
$x=im+j$ avec $0\le i,j<m$ (possible car $x<p-1\le m^2$). Alors
\[ g^{x}=A \iff g^{j}= A\,(g^{-m})^{i}. \]
\emph{Pas de bébé~:} on tabule les $m$ valeurs $g^{j}\bmod p$, $0\le j<m$, dans un
dictionnaire. \emph{Pas de géant~:} on calcule $u=g^{-m}\bmod p$ (par le petit théorème de
Fermat, $g^{-m}\equiv g^{\,p-1-m}$), puis les valeurs $A u^{i}$ pour $i=0,1,\dots$ jusqu'à en
trouver une dans la table~; on lit alors $j$ et l'on renvoie $x=im+j$. Coût~:
$O(\sqrt{p})$ opérations et $O(\sqrt{p})$ en espace.

\textbf{Ordres de grandeur pour $p$ de $2048$ bits.} Une machine à $10^{9}$ opérations par
seconde effectue environ $3{,}15\times10^{16}$ opérations par an.
\[ \text{naïf}~: p\approx 2^{2048}\approx 10^{616}\ \text{opérations},\qquad
   \text{BSGS}~: \sqrt{p}\approx 2^{1024}\approx 10^{308}\ \text{opérations}. \]
Même la variante en $\sqrt{p}$ demanderait de l'ordre de $10^{291}$ années — à comparer aux
$1{,}4\times 10^{10}$ années de l'âge de l'Univers. Et $O(\sqrt{p})$ en \emph{espace}, soit
$10^{308}$ entrées, est de toute façon physiquement impossible.

\textbf{Conclusion.} Le passage de $O(p)$ à $O(\sqrt{p})$ divise l'exposant par deux~: c'est
considérable en théorie, sans aucun effet pratique ici. C'est précisément la raison pour
laquelle la taille des clés est choisie de sorte que \emph{même} le meilleur algorithme connu
— le crible de corps de nombres, en complexité sous-exponentielle — reste hors de portée.
C'est aussi pourquoi on préfère aujourd'hui les courbes elliptiques~: aucun analogue du
crible n'y est connu, si bien qu'une clé de $256$ bits y offre le même niveau de sécurité
qu'une clé de $3072$ bits modulo $p$, pour un coût de calcul bien moindre.

\Q \textbf{Synthèse.}

\emph{Confidentialité persistante.} Si les exposants $x$ et $y$ sont éphémères et effacés
après usage, le secret $g^{xy}$ d'une session passée n'est reconstituable par personne~:
il ne subsiste nulle part, et la clé privée à long terme du serveur ne sert qu'à
\emph{signer} l'échange, jamais à le \emph{chiffrer}. Un attaquant qui aurait enregistré tout
le trafic puis obtiendrait la clé privée du serveur — par vol, contrainte légale ou faille —
ne pourrait toujours pas déchiffrer les sessions passées~; il pourrait seulement usurper
l'identité du serveur pour les sessions \emph{futures}. C'est la \emph{confidentialité
persistante} (\emph{forward secrecy}), rendue obligatoire par TLS~1.3.

Dans l'ancien mécanisme, au contraire, le client tirait la clé de session et l'envoyait
chiffrée avec la clé publique du serveur~: cette clé de session était donc, pour toujours,
récupérable par quiconque obtenait la clé privée. Une seule compromission ouvrait
rétroactivement des années d'archives — d'où sa suppression pure et simple.

\emph{Deux menaces de nature différente.} La partie~II traite de l'\textbf{intégrité} face à
un adversaire \emph{aveugle et non malveillant}~: le bruit, une mémoire défectueuse, un
routeur mal programmé. La somme de contrôle Internet est conçue contre les erreurs
\emph{aléatoires}, et l'on a vu qu'elle est très facile à contourner intentionnellement
(question 15~: il suffit de permuter deux mots). La partie~V traite de la
\textbf{confidentialité} et de l'\textbf{authenticité} face à un adversaire
\emph{intelligent}, qui choisit ses altérations pour ne pas être détecté.

Un mécanisme de la première catégorie ne protège en rien contre la seconde menace~: c'est
pourquoi TLS n'ajoute pas une somme de contrôle mais un \emph{code d'authentification de
message} muni d'une clé secrète, dont la sécurité repose sur une hypothèse calculatoire et
non sur une propriété combinatoire. Les deux couches coexistent parce qu'elles répondent à
des questions distinctes~: « ce paquet a-t-il été abîmé~? » et « ce paquet vient-il bien de
mon interlocuteur, et lui seul peut-il le lire~? ».

\vfill
\begin{center}\small\textsc{Fin du corrigé} \quad---\quad
\href{https://cpgeacademy.org}{cpgeacademy.org}\end{center}
\end{document}
